\chapter{Threats: danger zones}\label{ch:danger}

The map shows the threats each unit heard or named, as coloured zones. They also steer the routes in the next
chapter.

\section{Where threats come from}

\paragraph{Said by the soldier.} \code{danger.py} recognises six threat types in English and French: drone,
small-arms fire, sniper, explosions or shelling, armour, mines or IEDs. Negations (``no drones'') and injury
words (``gunshot to the chest'', ``blast'') don't count. A direction is taken from a clock hour, a compass word
(``from the north'', French \emph{est} only after a preposition), a numeric bearing (``sniper, bearing 1600
mils''), or ``overhead''. It can sit in the same clause or the next one (``taking fire, from my 3 o'clock'').

\paragraph{Heard by the microphone.} Gunfire (``Gunshot, gunfire'', ``Machine gun'', ``Fusillade'') and explosions
(``Explosion'', ``Artillery fire'', ``Boom'') come from YAMNet + AST: an alert fires if any 1 s window of YAMNet or
the whole clip in AST scores at least 0.3. Drones need a dedicated model, because AudioSet has no drone class. On
the phone or through a clean microphone we use an AST fine-tuned on drone audio (\code{preszzz} on Hugging Face,
86M parameters). Over the radio we use our own Near-small (Chapter~\ref{ch:models}). Wind on the microphone puts a
``doubtful'' flag on explosion alerts.

\section{From events to zones}

One phone can't locate a shooter or a drone. So each zone is centred on the \emph{reporter's} estimated position.
Its size is the distance at which that threat is usually heard or seen, plus the reporter's own uncertainty.

\begin{table}[htbp]\centering\small
\caption{Threat types (ranges and lifetimes are placeholders for the mentors).}\label{tab:threats}
\begin{tabular}{llrrl}
\toprule
Type & Colour & Range & Fades after & Packet code\\
\midrule
Drone & pink & 150 m & 10 min & U\\
Small-arms fire & red & 400 m & 15 min & G\\
Sniper & orange & 600 m & 30 min & S\\
Explosions / shelling & yellow & 800 m & 30 min & X\\
Armour & purple & 500 m & 20 min & A\\
Mines / IED & brown & 60 m & 24 h & M\\
\bottomrule
\end{tabular}
\end{table}

\begin{itemize}
  \item \textbf{Circle:} radius = position uncertainty + range (``overhead'': + half the range).
  \item \textbf{Wedge:} if a direction is given and the facing direction is known (heading confidence $\geq$ 5, from the
        sun, a compass, or walking direction), the zone becomes a 50$^\circ$ wedge of radius uncertainty + 1.5 $\times$ range.
  \item \textbf{Unplaced:} a reporter located worse than $\pm$500 m gets a ``?'' label instead of an area, and isn't used
        for routing.
  \item \textbf{Likely source area:} where zones from two or more units overlap (product of soft memberships on a 20 m
        grid), a thick ring marks it. It's only kept when the overlap is clearly smaller than the smallest zone.
  \item \textbf{Pinned in place:} a unit that walks away leaves its threats where it reported them.
  \item \textbf{``Near you'':} a unit standing inside another unit's zone is warned in its read-back.
\end{itemize}

\nanofig{danger}{Danger zones on the coordinator map: circles and wedges around reporters, dashed when heard by a microphone,
and a thick ring where two units' reports overlap.}
{Top-down dark tactical map of a small town, 16:9, the map drawn in near-black navy tones as a rounded rectangle on a plain white background, thin grey streets and building blocks.
A translucent red circle with a solid outline labelled ``SMALL-ARMS FIRE · voice'' around a unit marker ``C1''. A
second translucent red circle with a dashed outline labelled ``SMALL-ARMS FIRE · heard'' around unit ``C2''. Where the
two circles overlap, a thick red ring labelled ``likely source area''. A pink dashed circle labelled ``DRONE'' around unit
``C3''. An orange 50-degree wedge labelled ``SNIPER · bearing 1600 mils'' pointing east from unit ``C4''. A small legend
box at bottom right with coloured squares: Drone, Small-arms fire, Sniper, Explosions, Armour, Mines. Exactly these labels. Colour palette (use only these): navy \#0D1110, cyan \#1AA7C4, amber \#E09A2B, red \#E5484D, green \#3FA877, grey \#8B9A94, plus white for text on dark shapes.}

\section{Results on real recordings}

\begin{table}[htbp]\centering\small
\caption{Danger detection on real recordings, laptop models, clean audio (\code{scripts/eval\_danger.py}).}\label{tab:dangereval}
\begin{tabularx}{\linewidth}{Xl}
\toprule
Test & Result\\
\midrule
88 real gunshot recordings (HEAR gunshot triangulation) & small-arms alert on 97 \% (small arms or explosion: 100 \%)\\
150 drone clips (DADS) & drone alert on 72 \% (first evaluation; 91 \% on another 150-clip shard)\\
85 no-danger ESC-50 clips (helicopter, airplane, engine, door knocks, wind...) & false alarm on 2 \% (wind read as explosion, now flagged doubtful)\\
\bottomrule
\end{tabularx}
\end{table}

The drone detector's model card claims 99.6 \%, which we don't reproduce. It may also have been trained on DADS, so
our drone recall is optimistic. The more useful number is zero drone false alarms on helicopters and airplanes.
